\documentclass[10pt]{article}
\usepackage[utf8]{inputenc}
\usepackage[T1]{fontenc}
\usepackage{lmodern}
\usepackage{enumitem}
\usepackage{multicol}
\usepackage{geometry}
\usepackage{setspace}
\geometry{a4paper, margin=0.75in}
\setstretch{0.8}
\begin{document}

\section*{IoT Systems and Technologies \footnotesize [B99CB8CB1]}
{\footnotesize Nome:\_\_\_\_\_\_\_\_\_\_\_\_\_\_\_\_ Cognome:\_\_\_\_\_\_\_\_\_\_\_\_\_\_\_\_ Mat:\_\_\_\_\_\_\_\_}
\begin{multicols}{2}

\textbf{Q1: Cosa rende un sistema embedded un dispositivo IoT?}

\begin{enumerate}[label=\arabic*.]
\item La connessione in rete
\item L’integrazione con un PC
\item L’uso di un'app mobile
\item La presenza di una batteria
\end{enumerate}
\textbf{Q2: Quale elemento riduce la probabilità di meta-stabilità?}

\begin{enumerate}[label=\arabic*.]
\item EEPROM
\item DAC
\item Schmitt-trigger
\item ADC
\end{enumerate}
\textbf{Q3: Cos'è la meta-stabilità?}

\begin{enumerate}[label=\arabic*.]
\item Stato incerto di un latch per segnale indeterminato
\item Valore massimo di tensione
\item Stato di reset del microcontrollore
\item Errore di memoria
\end{enumerate}
\textbf{Q4: In una interfaccia single-ended}

\begin{enumerate}[label=\arabic*.]
\item E' meno soggetta ad interferenze e rumore
\item Il trasmettitore ed il ricevitore fanno riferimento alla stessa massa
\item Ogni bit trasmesso è seguito da un bit di end
\item Ogni bit è trasmesso su due linee singole (cioè non comunicanti tra loro)
\end{enumerate}
\textbf{Q5: Come viene rilevato l’inizio di un pacchetto UART?}

\begin{enumerate}[label=\arabic*.]
\item Dall’assenza di segnale
\item Dal fronte di discesa del bit di start
\item Dal bit di parità
\item Dal clock condiviso
\end{enumerate}
\textbf{Q6: In una interfaccia UART:}

\begin{enumerate}[label=\arabic*.]
\item Utilizza un bus con le linee RX e TX
\item Le linee di trasmissione sono connesse tra di loro
\item La linea di trasmissione è connessa a quella di ricezione
\item Nessuna delle opzioni riportate
\end{enumerate}
\textbf{Q7: Per un LED, tipicamente è necessaria:}

\begin{enumerate}[label=\arabic*.]
\item Resistenza di qualche centinaio di Ohm in parallelo
\item Resistenza di qualche Mega Ohm in serie
\item Resistenza di qualche Mega Ohm in parallelo
\item Resistenza di qualche centinaio di Ohm in serie
\end{enumerate}
\textbf{Q8: Se connetto una board ESP32 all'alimentazione:}

\begin{enumerate}[label=\arabic*.]
\item Nessuna delle altre risposte
\item Inizia ad eseguire l'ultimo codice caricato sulla flash
\item Non esegue niente finchè non si immette un nuovo codice sulla flash
\item Azzera il proprio contenuto flash
\end{enumerate}
\textbf{Q9: Cosa accade se si preme il bottone durante un `delay()`?}

\begin{enumerate}[label=\arabic*.]
\item Il microcontrollore non se ne accorge
\item Il segnale del pin è messo in coda alle richieste
\item La resistenza interna di pull-up cambia valore
\item L’ESP32 si resetta
\end{enumerate}
\textbf{Q10: Quando si usa dacWrite invece di PWM?}

\begin{enumerate}[label=\arabic*.]
\item Quando si usano motori
\item Quando si vuole risparmiare memoria
\item Quando serve un segnale digitale
\item Quando serve un’uscita analogica precisa
\end{enumerate}
\textbf{Q11: Perché il debounce è importante con gli interrupt?}

\begin{enumerate}[label=\arabic*.]
\item Per evitare chiamate multiple della ISR
\item Per aumentare la luminosità del LED
\item Per risparmiare energia
\item Per evitare la disconnessione del pin
\end{enumerate}
\textbf{Q12: Cosa rappresenta un "Device Template" in Blynk?}

\begin{enumerate}[label=\arabic*.]
\item Un'app per smartphone per progettare interfacce IoT
\item Un modello per progettare sensori
\item Un modello per creare dispositivi con le stesse funzionalità
\item Un firmware aggiornato per modelli con simili funzionalità
\end{enumerate}
\textbf{Q13: Dato una gestione dell’I/O basata su DMA per trasferire burst di dati, si assuma che il periodo di silenzio  ($T_{off}$) di ciascun burst venga utilizzato per programmare il DMA e per gestire l’interrupt di fine trasferimento e che nel periodo $T_{ON} = 0.2$ sec arrivino un massimo di $N=10^6$ bytes. Calcolare il numero massimo $N_{max}$ di bytes in ciascun burst che possono essere trasferiti attraverso il DMA in una architettura Harvard con percentuale di conflitti $f = 0.3$ gestiti con DMA waits for CPU e dato $T_{LOAD\/STORE} = 2 \times T_{tr\_DMA}$ e $T_{tr\_DMA} = 10^{-6}$ sec/byte.}

\begin{enumerate}[label=\arabic*.]
\item $0.25 \times 10^5$ bytes
\item $0.71 \times 2 \times 10^6$ bytes
\item $0.125 \times 10^6$ bytes
\item $2 \times 10^6$ bytes
\end{enumerate}
\textbf{Q14: In un sistema IoT occorre gestire l'acquisizione di burst di dati caratterizzati da un tempo di silenzio $t_{off}=50\mu sec$. Ogni burst ha una durata $t_{ON}$ che puo' variare da $t_{ON\_min}$ uguale a $0.1 sec$ a $4 sec$. In ciascun periodo $t_{ON}$ arrivano sempre 2MB, con arrivi distribuiti uniformemente.  Assumendo una gestione con interrupt e overhead $W=0.3$ calcolare il valore massimo di $T_{ISR}$ per acquisire correttamente i dati del traffico bursty sopra descritto.}

\begin{enumerate}[label=\arabic*.]
\item $25\times10^{-5}$ sec/byte
\item $0.01\times10^{-6}$ sec/byte
\item $3.8\times10^{-8}$ sec/byte
\item Nessuna delle altre opzioni
\end{enumerate}
\textbf{Q15: Assumendo una gestion del sistema di I/O con polling, con $T_{TR}=100\mu sec$, calcolare il minimo valore di $\eta_{Poll}$ affinche' si possano acquisire 5 sorgenti analogiche ciascuna caratterizzata da una $f_{MAX[segnale]}$ di 0.3Kbyte/sec.  Indicare quale dei seguenti valore è quello corretto:}

\begin{enumerate}[label=\arabic*.]
\item 0.35
\item 0.15
\item 0.3
\item 0.27
\end{enumerate}
\textbf{Q16: In un sistema di I/O gestito con il DMA, architettura Harvard, $f=0.25$, il max throughput di I/O ottenibile scambiando N bytes:}

\begin{enumerate}[label=\arabic*.]
\item E' inversamente proporzionale al valore N
\item E' direttamente proporzionale al valore N
\item Nessuna delle altre opzioni
\item E' almeno il doppio di quello ottenibile con $f=0$
\end{enumerate}
\textbf{Q17: Si utilizzi un sistema di gestione dell’I/O basato sul DMA per trasferire N bytes di dati dall’I/O. Si assuma una architettura Harvard, una percentuale di conflitti f diversa da 0 e una politica "CPU waits for DMA". Quale fra le seguenti affermazioni è corretta:}

\begin{enumerate}[label=\arabic*.]
\item Il Throughput di I/O diminuisce all’aumentare di f
\item Nessuna delle altre opzioni
\item Il Throughput di I/O raddoppia se f si dimezza
\item Il Throughput di I/O dipende dal valore di f
\end{enumerate}
\textbf{Q18: In un sistema IoT con due sorgenti analogiche caratterizzate dallo stesso valore di $f_{MAX[segnale]}$ espresso in byte/sec, assumendo una catena di acquisizione con un unico convertitore A/D, selezionare quale delle seguenti opzioni è vera:}

\begin{enumerate}[label=\arabic*.]
\item Il massimo throughput di I/O deve essere almeno pari a $4 \times f_{MAX[segnale]}$ byte/sec
\item Nessuna delle opzioni
\item Il massimo throughput di I/O deve essere minore o uguale a $8 \times f_{MAX[segnale]}$ byte/sec
\item Il massimo throughput di I/O deve essere minore o uguale a $2 \times f_{MAX[segnale]}$ byte/sec
\end{enumerate}
\textbf{Q19: Considerando una gestione dell’I/O con interrupt:}

\begin{enumerate}[label=\arabic*.]
\item Il vettore di interrupt viene generato dalla CPU
\item Il vettore di interrupt elimina il problema di gestire la priorità
\item Nessuna delle opzioni riportate
\item L’utilizzo della tecnica dell’interrupt vettorizzato non diminuisce l’interrupt overhead W
\end{enumerate}
\textbf{Q20: In un sistema di I/O gestito con il polling, la massima frequenza di segnale $f_{MAX[segnale]}$ acquisibile:}

\begin{enumerate}[label=\arabic*.]
\item Nessuna delle altre opzioni
\item Non dipende da W
\item E' inversamente proporzionale al valore $T_{TR}$
\item Non dipende da $X_M$
\end{enumerate}
\textbf{Q21: Con riferimento all’architettura generale di un sistema IoT dire quale fra le seguenti opzioni è quella vera}

\begin{enumerate}[label=\arabic*.]
\item il Fog Computing serve soprattutto a esporre interfacce di programmazione più agevoli agli sviluppatori di applicazioni, nascondendo loro la complessità dei livelli più bassi della gerarchia
\item Il Fog Computing non si utilizza nel caso vi siano vincoli troppo stringenti di real time
\item Nessuna delle altre risposte
\item Il Fog Computing serve soprattutto a interfacciare i nodi IoT con la piattaforma IoT
\end{enumerate}
\end{multicols}

\newpage

\section*{IoT Systems and Technologies \footnotesize [B99CB8CB2]}
{\footnotesize Nome:\_\_\_\_\_\_\_\_\_\_\_\_\_\_\_\_ Cognome:\_\_\_\_\_\_\_\_\_\_\_\_\_\_\_\_ Mat:\_\_\_\_\_\_\_\_}
\begin{multicols}{2}

\textbf{Q1: Qual è la relazione tra Embedded System e IoT?}

\begin{enumerate}[label=\arabic*.]
\item L’IoT è una parte del sistema embedded
\item Sono la stessa cosa
\item Un IoT device contiene un sistema embedded
\item Il sistema embedded è solo software
\end{enumerate}
\textbf{Q2: In un microcontrollore, l'output PWM di un pin può essere utilizzato}

\begin{enumerate}[label=\arabic*.]
\item Per programmare le Interrupt Service Routine
\item Per generare un segnale di clock interno a risoluzione maggiore
\item Per estendere il range di valori di un timer convenzionale
\item Per implementare un convertitore da digitale ad analogico
\end{enumerate}
\textbf{Q3: Qual è lo svantaggio del circuito resistivo binary-weight per DAC?}

\begin{enumerate}[label=\arabic*.]
\item Richiede resistenze di precisione
\item Ha alta latenza
\item Non funziona con segnali digitali
\item È troppo lento
\end{enumerate}
\textbf{Q4: Un’interfaccia single-ended misura...}

\begin{enumerate}[label=\arabic*.]
\item Solo impulsi digitali
\item La frequenza assoluta
\item La differenza tra due segnali
\item La tensione rispetto al GND comune
\end{enumerate}
\textbf{Q5: L'interfaccia SPI è...}

\begin{enumerate}[label=\arabic*.]
\item Differenziale e half-duplex
\item Sincrona, punto a punto, master-slave
\item Asincrona, broadcast
\item Parallela e half-duplex
\end{enumerate}
\textbf{Q6: Come viene rilevato l’inizio di un pacchetto UART?}

\begin{enumerate}[label=\arabic*.]
\item Dall’assenza di segnale
\item Dal fronte di discesa del bit di start
\item Dal bit di parità
\item Dal clock condiviso
\end{enumerate}
\textbf{Q7: A cosa serve la fase di calibrazione nei sensori luce?}

\begin{enumerate}[label=\arabic*.]
\item Per impostare il colore del LED
\item Per impostare la frequenza del LED
\item Per stabilire un valore fisso di tensione
\item Per rilevare il range di valori utili
\end{enumerate}
\textbf{Q8: La calibrazione di un sensore potrebbe essere necessaria:}

\begin{enumerate}[label=\arabic*.]
\item Per determinare il range di variazione delle tensioni lette
\item Per determinare la corrente massima assorbita
\item Per riavviare il circuito del sensore in caso di inutilizzo prolungato
\item Per evitare che la DigitalRead restituisca valori intermedi tra 0 ed 1
\end{enumerate}
\textbf{Q9: Quando si usa dacWrite invece di PWM?}

\begin{enumerate}[label=\arabic*.]
\item Quando serve un segnale digitale
\item Quando si vuole risparmiare memoria
\item Quando si usano motori
\item Quando serve un’uscita analogica precisa
\end{enumerate}
\textbf{Q10: Cosa accade se si preme il bottone durante un `delay()`?}

\begin{enumerate}[label=\arabic*.]
\item La resistenza interna di pull-up cambia valore
\item L’ESP32 si resetta
\item Il microcontrollore non se ne accorge
\item Il segnale del pin è messo in coda alle richieste
\end{enumerate}
\textbf{Q11: Cosa è inopportuno porre dentro una ISR:}

\begin{enumerate}[label=\arabic*.]
\item Un insieme di istruzioni che usi variabili booleane
\item Un insieme di istruzioni che modificano variabili globali
\item Un insieme di istruzioni che modificano variabili locali
\item Un insieme di istruzioni ad alta latenza
\end{enumerate}
\textbf{Q12: A cosa serve $BlynkTimer$?}

\begin{enumerate}[label=\arabic*.]
\item Ad evitare di usare $delay()$
\item Ad aumentare la potenza del segnale Wi-Fi per il cloud
\item Per modificare i pin virtuali per il cloud
\item A controllare il drift del clock della CPU
\end{enumerate}
\textbf{Q13: Dato una gestione dell’I/O basata su DMA per trasferire burst di dati, si assuma che il periodo di silenzio  ($T_{off}$) di ciascun burst venga utilizzato per programmare il DMA e per gestire l’interrupt di fine trasferimento e che nel periodo $T_{ON} = 0.2$ sec arrivino un massimo di $N=10^6$ bytes. Calcolare il numero massimo $N_{max}$ di bytes in ciascun burst che possono essere trasferiti attraverso il DMA in una architettura Harvard con percentuale di conflitti $f = 0.3$ gestiti con DMA waits for CPU e dato $T_{LOAD\/STORE} = 2 \times T_{tr\_DMA}$ e $T_{tr\_DMA} = 10^{-6}$ sec/byte.}

\begin{enumerate}[label=\arabic*.]
\item $0.25 \times 10^5$ bytes
\item $0.71 \times 2 \times 10^6$ bytes
\item $0.125 \times 10^6$ bytes
\item $2 \times 10^6$ bytes
\end{enumerate}
\textbf{Q14: Assumendo una gestione del sistema di I/O con Interrupt, con $T_{ISR}=100\mu sec$, calcolare il valore massimo di $W$ affinche' si possano acquisire 10 sorgenti analogiche ciascuna caratterizzata da una $f_{MAX[segnale]}$ di 0.3Kbyte/sec.}

\begin{enumerate}[label=\arabic*.]
\item 0.66
\item 0.73
\item 0.5
\item 0.81
\end{enumerate}
\textbf{Q15: Assumendo una gestion del sistema di I/O con polling, con $T_{TR}=100\mu sec$, calcolare il minimo valore di $\eta_{Poll}$ affinche' si possano acquisire 5 sorgenti analogiche ciascuna caratterizzata da una $f_{MAX[segnale]}$ di 0.3Kbyte/sec.  Indicare quale dei seguenti valore è quello corretto:}

\begin{enumerate}[label=\arabic*.]
\item 0.35
\item 0.27
\item 0.15
\item 0.3
\end{enumerate}
\textbf{Q16: Si utilizzi un sistema di gestione dell’I/O basato sul DMA per trasferire N bytes di dati dall’I/O. Si assuma una architettura Harvard, una percentuale di conflitti f diversa da 0 e una politica "CPU waits for DMA". Quale fra le seguenti affermazioni è corretta:}

\begin{enumerate}[label=\arabic*.]
\item Nessuna delle altre opzioni
\item Il $CPU_{TIME\_IO}$ che la CPU spende per l’I/O dipende dal prodotto $f \times N$
\item Il $CPU_{TIME\_IO}$ che la CPU spende per l’I/O non dipende da f
\item Il $CPU_{TIME\_IO}$ che la CPU spende per l’I/O è trascurabile per valori di N molto grandi
\end{enumerate}
\textbf{Q17: Data una gestione dell’I/O basata su DMA per trasferire burst di dati dall’I/O e che un periodo di silenzio $T_{off}$ di ciascun burst venga utilizzato per programmare il DMA e per gestire l’interrupt di fine trasferimento. Qual è il massimo throughput che si può smaltire in ciascun burst, nel periodo $T_{ON}$ in cui arrivano N bytes, nel caso di una architettura Harvard, percentuale di conflitti f non nulla, e una politica di gestione dei conflitti di tipo "DMA waiting for CPU" ?}

\begin{enumerate}[label=\arabic*.]
\item $\frac{N}{s+ 2 \times f \times N + N} \times \frac{1}{T_{tr\_DMA}} $
\item $\frac{1}{T_{tr\_DMA}}$
\item $0.5 \times \frac{N}{s+ 2 \times f \times N + N} \times \frac{1}{T_{tr\_DMA}} $
\item Nessuna delle altre opzioni
\end{enumerate}
\textbf{Q18: Con riferimento alla catena di acquisizione di un sistema IoT con una sorgente analogica caratterizzata da un valore di $f_{MAX[segnale]} = K$, selezionare quale delle seguenti opzioni è vera:}

\begin{enumerate}[label=\arabic*.]
\item La latenza del convertitore A/D deve essere almeno il doppio di $f_{MAX[segnale]}$
\item Il throughput minimo del sistema di I/O deve essere almeno il doppio della frequenza di campionamento di Nyquist
\item La latenza del Sample\&Hold deve essere il doppio di $f_{MAX[segnale]}$
\item Nessuna delle altre opzioni
\end{enumerate}
\textbf{Q19: Considerando l’utilizzo dell’interrupt come tecnica di gestione dell’I/O:}

\begin{enumerate}[label=\arabic*.]
\item Il vettore di interrupt serve principalmente ad indicare alla CPU la priorità di gestione della ISR
\item Nessuna delle opzioni riportate
\item La gestione della priorità serve ad abbattere l’interrupt overhead
\item L’interrupt overhead W può essere diminuito se si utilizza una gestione vettorizzata dell’interrupt
\end{enumerate}
\textbf{Q20: In un sistema di I/O gestito con il polling, la massima frequenza di segnale $f_{MAX[segnale]}$ acquisibile:}

\begin{enumerate}[label=\arabic*.]
\item E' direttamente proporzionale al valore $T_{TR}$
\item Non dipende da $X_M$
\item Nessuna delle altre opzioni
\item Non dipende da W
\end{enumerate}
\textbf{Q21: Con riferimento all’architettura generale di un sistema IoT dire quale fra le seguenti opzioni è quella vera}

\begin{enumerate}[label=\arabic*.]
\item Nessuna delle altre risposte
\item L’uso del MIST Computing, in alcuni scenari che prevedono comunicazioni wireless, comporta un maggior consumo di energia del nodo IoT dovuto alla trasmissione dati
\item Il ruolo del MIST Computing è principalmente quello di ridurre la complessità del software a bordo dei nodi IoT per limitare il consumo di potenza del nodo
\item L’uso congiunto di Fog e MIST computing in un sistema IoT normalmente serve ad evitare l’utilizzo del Cloud Computing
\end{enumerate}
\end{multicols}

\newpage

\section*{IoT Systems and Technologies \footnotesize [B99CB8CB3]}
{\footnotesize Nome:\_\_\_\_\_\_\_\_\_\_\_\_\_\_\_\_ Cognome:\_\_\_\_\_\_\_\_\_\_\_\_\_\_\_\_ Mat:\_\_\_\_\_\_\_\_}
\begin{multicols}{2}

\textbf{Q1: Cosa contiene la memoria flash di un microcontrollore?}

\begin{enumerate}[label=\arabic*.]
\item Le connessioni di rete attive
\item Il programma da eseguire
\item I dati per il debug mediante connessione seriale
\item I dati temporanei delle variabili
\end{enumerate}
\textbf{Q2: Qual è uno svantaggio dei sincronizzatori?}

\begin{enumerate}[label=\arabic*.]
\item Introducono un ritardo nella trasmissione del segnale
\item Richiedono più memoria
\item Aumentano il rumore
\item Consumano più corrente
\end{enumerate}
\textbf{Q3: Un pin che supporta il PWM con segnali di 0 e 5V e con un periodo di 1 millisecondo:}

\begin{enumerate}[label=\arabic*.]
\item Non potrebbe generare segnali di durata  minore di un 1 millisecondo
\item Può generare qualunque valore di output che sia compreso tra i 0 ed i 5V
\item Genera un segnale di circa 0.2V
\item Potrebbe generare un segnale di circa 1V restando al valore alto per 200 microsecondi nel periodo
\end{enumerate}
\textbf{Q4: Nella comunicazione sincrona...}

\begin{enumerate}[label=\arabic*.]
\item Si usano sempre 4 linee dati
\item Ogni dispositivo usa il proprio clock
\item I clock di trasmissione e ricezione sono collegati
\item Non è necessaria una linea di clock
\end{enumerate}
\textbf{Q5: Come viene rilevato l’inizio di un pacchetto UART?}

\begin{enumerate}[label=\arabic*.]
\item Dal fronte di discesa del bit di start
\item Dal clock condiviso
\item Dall’assenza di segnale
\item Dal bit di parità
\end{enumerate}
\textbf{Q6: Quali sono le quattro linee dell’interfaccia SPI?}

\begin{enumerate}[label=\arabic*.]
\item MOSI, MISO, SCK, SS
\item IN, OUT, GND, VCC
\item TXD, RXD, CLK, EN
\item SDA, SCL, ACK, RST
\end{enumerate}
\textbf{Q7: Per un LED, tipicamente è necessaria:}

\begin{enumerate}[label=\arabic*.]
\item Resistenza di qualche Mega Ohm in serie
\item Resistenza di qualche Mega Ohm in parallelo
\item Resistenza di qualche centinaio di Ohm in serie
\item Resistenza di qualche centinaio di Ohm in parallelo
\end{enumerate}
\textbf{Q8: Qual è il problema del “floating pin” quando un pulsante è aperto?}

\begin{enumerate}[label=\arabic*.]
\item Il pin si danneggia
\item Il pin può assumere valori casuali
\item Il LED si spegne automaticamente
\item La corrente diventa negativa
\end{enumerate}
\textbf{Q9: Cosa accade se si preme il bottone durante un `delay()`?}

\begin{enumerate}[label=\arabic*.]
\item Il segnale del pin è messo in coda alle richieste
\item La resistenza interna di pull-up cambia valore
\item L’ESP32 si resetta
\item Il microcontrollore non se ne accorge
\end{enumerate}
\textbf{Q10: Nel codice di ESP32 la funzione setup()}

\begin{enumerate}[label=\arabic*.]
\item Viene eseguita una sola volta, all'uscita della funzione loop()
\item Può essere eseguita in qualunque momento, ma una volta sola, mediante apposita chiamata
\item Viene eseguita alla fine di ogni ripetizione della funzione loop()
\item Viene eseguita una sola volta, ad ogni avvio del programma
\end{enumerate}
\textbf{Q11: Perché il debounce è importante con gli interrupt?}

\begin{enumerate}[label=\arabic*.]
\item Per evitare la disconnessione del pin
\item Per evitare chiamate multiple della ISR
\item Per aumentare la luminosità del LED
\item Per risparmiare energia
\end{enumerate}
\textbf{Q12: A cose serve la funzione $Blink.run()$}

\begin{enumerate}[label=\arabic*.]
\item Apre una nuova connessione Wi-Fi per gli interrupt
\item Gestisce gli interrupt in alternativa ad attachInterrupt()
\item Configura i pin GPIO per il cloud dentro setup()
\item Gestisce la connessione cloud e messaggi virtuali dentro loop()
\end{enumerate}
\textbf{Q13: Dato una gestione dell’I/O basata su DMA per trasferire burst di dati, si assuma che il periodo di silenzio  ($T_{off}$) di ciascun burst venga utilizzato per programmare il DMA e per gestire l’interrupt di fine trasferimento e che nel periodo $T_{ON} = 0.2$ sec arrivino un massimo di $N=10^6$ bytes. Calcolare il numero massimo $N_{max}$ di bytes in ciascun burst che possono essere trasferiti attraverso il DMA in una architettura Harvard con percentuale di conflitti $f = 0.3$ gestiti con DMA waits for CPU e dato $T_{LOAD\/STORE} = 2 \times T_{tr\_DMA}$ e $T_{tr\_DMA} = 10^{-6}$ sec/byte.}

\begin{enumerate}[label=\arabic*.]
\item $2 \times 10^6$ bytes
\item $0.71 \times 2 \times 10^6$ bytes
\item $0.25 \times 10^5$ bytes
\item $0.125 \times 10^6$ bytes
\end{enumerate}
\textbf{Q14: Sia dato un sistema IoT in cui si voglia valutare la massima frequenza di segnale acquisibile da una sorgente analogica, espressa in byte/sec. Assumendo una gestione del sistema di I/O con Interrupt, con $T_{ISR}=100\mu sec$ e $W=0.1$, indicare la risposta corretta:}

\begin{enumerate}[label=\arabic*.]
\item $0.45\times10^4 byte/sec$
\item $0.22\times10^5 byte/sec$
\item $0.045\times10^3 byte/sec$
\item $0.32\times10^3 byte/sec$
\end{enumerate}
\textbf{Q15: Assumendo una gestion del sistema di I/O con polling, con $T_{TR}=100\mu sec$, calcolare il minimo valore di $\eta_{Poll}$ affinche' si possano acquisire 5 sorgenti analogiche ciascuna caratterizzata da una $f_{MAX[segnale]}$ di 0.3Kbyte/sec.  Indicare quale dei seguenti valore è quello corretto:}

\begin{enumerate}[label=\arabic*.]
\item 0.35
\item 0.3
\item 0.15
\item 0.27
\end{enumerate}
\textbf{Q16: In un sistema di I/O gestito con il DMA, architettura Harvard, $f=0.25$, il max throughput di I/O ottenibile scambiando N bytes:}

\begin{enumerate}[label=\arabic*.]
\item E' direttamente proporzionale al valore N
\item E' inversamente proporzionale al valore N
\item E' almeno il doppio di quello ottenibile con $f=0$
\item Nessuna delle altre opzioni
\end{enumerate}
\textbf{Q17: Si utilizzi un sistema di gestione dell’I/O basato sul DMA per trasferire N bytes di dati dall’I/O. Si assuma una architettura Harvard, una percentuale di conflitti f diversa da 0 e una politica "CPU waits for DMA". Quale fra le seguenti affermazioni è corretta:}

\begin{enumerate}[label=\arabic*.]
\item Il Throughput di I/O dipende dal valore di f
\item Nessuna delle altre opzioni
\item Il Throughput di I/O diminuisce all’aumentare di f
\item Il Throughput di I/O raddoppia se f si dimezza
\end{enumerate}
\textbf{Q18: In un sistema IoT con due sorgenti analogiche caratterizzate dallo stesso valore di $f_{MAX[segnale]}$ espresso in byte/sec, assumendo una catena di acquisizione con un unico convertitore A/D, selezionare quale delle seguenti opzioni è vera:}

\begin{enumerate}[label=\arabic*.]
\item Il massimo throughput di I/O deve essere minore o uguale a $2 \times f_{MAX[segnale]}$ byte/sec
\item Nessuna delle opzioni
\item Il massimo throughput di I/O deve essere minore o uguale a $8 \times f_{MAX[segnale]}$ byte/sec
\item Il massimo throughput di I/O deve essere almeno pari a $4 \times f_{MAX[segnale]}$ byte/sec
\end{enumerate}
\textbf{Q19: Considerando una gestione dell’I/O con interrupt:}

\begin{enumerate}[label=\arabic*.]
\item L’utilizzo della tecnica dell’interrupt vettorizzato non diminuisce l’interrupt overhead W
\item Il vettore di interrupt elimina il problema di gestire la priorità
\item Nessuna delle opzioni riportate
\item Il vettore di interrupt viene generato dalla CPU
\end{enumerate}
\textbf{Q20: In un sistema di I/O gestito con il polling, la massima frequenza di segnale $f_{MAX[segnale]}$ acquisibile:}

\begin{enumerate}[label=\arabic*.]
\item Non dipende da W
\item Nessuna delle altre opzioni
\item Non dipende da $X_M$
\item E' direttamente proporzionale al valore $T_{TR}$
\end{enumerate}
\textbf{Q21: Con riferimento all’architettura generale di un sistema IoT dire quale fra le seguenti opzioni è quella vera}

\begin{enumerate}[label=\arabic*.]
\item il Fog Computing serve soprattutto a esporre interfacce di programmazione più agevoli agli sviluppatori di applicazioni, nascondendo loro la complessità dei livelli più bassi della gerarchia
\item Nessuna delle altre risposte
\item Il Fog Computing non si utilizza nel caso vi siano vincoli troppo stringenti di real time
\item Il Fog Computing serve soprattutto a interfacciare i nodi IoT con la piattaforma IoT
\end{enumerate}
\end{multicols}

\newpage

\section*{IoT Systems and Technologies \footnotesize [B99CB8CB4]}
{\footnotesize Nome:\_\_\_\_\_\_\_\_\_\_\_\_\_\_\_\_ Cognome:\_\_\_\_\_\_\_\_\_\_\_\_\_\_\_\_ Mat:\_\_\_\_\_\_\_\_}
\begin{multicols}{2}

\textbf{Q1: Cosa contiene la memoria flash di un microcontrollore?}

\begin{enumerate}[label=\arabic*.]
\item I dati per il debug mediante connessione seriale
\item Le connessioni di rete attive
\item I dati temporanei delle variabili
\item Il programma da eseguire
\end{enumerate}
\textbf{Q2: Quale elemento riduce la probabilità di meta-stabilità?}

\begin{enumerate}[label=\arabic*.]
\item ADC
\item DAC
\item Schmitt-trigger
\item EEPROM
\end{enumerate}
\textbf{Q3: Qual è uno svantaggio del tracking converter?}

\begin{enumerate}[label=\arabic*.]
\item È lento a seguire variazioni rapide del segnale
\item Non può usare PWM
\item Ha solo 2 bit di risoluzione
\item Richiede molta RAM
\end{enumerate}
\textbf{Q4: Nella comunicazione sincrona...}

\begin{enumerate}[label=\arabic*.]
\item Si usano sempre 4 linee dati
\item I clock di trasmissione e ricezione sono collegati
\item Ogni dispositivo usa il proprio clock
\item Non è necessaria una linea di clock
\end{enumerate}
\textbf{Q5: Come viene gestita la dominanza in I2C?}

\begin{enumerate}[label=\arabic*.]
\item Il livello alto ha priorità
\item Il livello basso vince sempre
\item Dipende dal master
\item Si alternano i nodi
\end{enumerate}
\textbf{Q6: In una interfaccia I2C:}

\begin{enumerate}[label=\arabic*.]
\item Seleziona uno slave mediante pilotando il segnale di clock
\item Presenta un bus con due linee delle MOSI e MISO
\item Il bus è costituito da due linee, di cui una è il clock
\item Il bus è costituito da due linee, più una terza linea per il clock
\end{enumerate}
\textbf{Q7: I sensori a 2 pin funzionano tipicamente come:}

\begin{enumerate}[label=\arabic*.]
\item Divisori di tensione
\item Amplificatori operazionali
\item DAC digitali
\item Filtri attivi
\end{enumerate}
\textbf{Q8: Qual è una soluzione al problema del bouncing?}

\begin{enumerate}[label=\arabic*.]
\item Cambiare il pin GPIO
\item Ignorare interrupt se è trascorso poco tempo dall'ultima chiamata
\item Utilizzare un LED come filtro
\item Disattivare il Wi-Fi e le connessioni seriali
\end{enumerate}
\textbf{Q9: Qual è la funzione di `ledcAttach()`?}

\begin{enumerate}[label=\arabic*.]
\item Configurare un pin analogico
\item Attaccare un interrupt
\item Stabilire comunicazione UART
\item Configurare il PWM avanzato
\end{enumerate}
\textbf{Q10: Su quali pin è possibile utilizzare `dacWrite()`?}

\begin{enumerate}[label=\arabic*.]
\item Solo su alcuni pin
\item Su pin che vanno scelti in fase di setup
\item Su tutti i pin digitali
\item Su tutti i pin analogici
\end{enumerate}
\textbf{Q11: Cosa è inopportuno porre dentro una ISR:}

\begin{enumerate}[label=\arabic*.]
\item Un insieme di istruzioni che modificano variabili locali
\item Un insieme di istruzioni che modificano variabili globali
\item Un insieme di istruzioni che usi variabili booleane
\item Un insieme di istruzioni ad alta latenza
\end{enumerate}
\textbf{Q12: A cose serve la funzione $Blink.run()$}

\begin{enumerate}[label=\arabic*.]
\item Gestisce gli interrupt in alternativa ad attachInterrupt()
\item Configura i pin GPIO per il cloud dentro setup()
\item Gestisce la connessione cloud e messaggi virtuali dentro loop()
\item Apre una nuova connessione Wi-Fi per gli interrupt
\end{enumerate}
\textbf{Q13: Dato una gestione dell’I/O basata su DMA per trasferire burst di dati, si assuma che il periodo di silenzio  ($T_{off}$) di ciascun burst venga utilizzato per programmare il DMA e per gestire l’interrupt di fine trasferimento e che nel periodo $T_{ON} = 0.2$ sec arrivino un massimo di $N=10^6$ bytes. Calcolare il numero massimo $N_{max}$ di bytes in ciascun burst che possono essere trasferiti attraverso il DMA in una architettura Harvard con percentuale di conflitti $f = 0.3$ gestiti con DMA waits for CPU e dato $T_{LOAD\/STORE} = 2 \times T_{tr\_DMA}$ e $T_{tr\_DMA} = 10^{-6}$ sec/byte.}

\begin{enumerate}[label=\arabic*.]
\item $0.25 \times 10^5$ bytes
\item $2 \times 10^6$ bytes
\item $0.125 \times 10^6$ bytes
\item $0.71 \times 2 \times 10^6$ bytes
\end{enumerate}
\textbf{Q14: Assumendo una gestione del sistema di I/O con Interrupt, con $T_{ISR}=100\mu sec$, calcolare il valore massimo di $W$ affinche' si possano acquisire 10 sorgenti analogiche ciascuna caratterizzata da una $f_{MAX[segnale]}$ di 0.3Kbyte/sec.}

\begin{enumerate}[label=\arabic*.]
\item 0.66
\item 0.81
\item 0.73
\item 0.5
\end{enumerate}
\textbf{Q15: In un sistema IoT occorre gestire l'acquisizione di un burst di dati con tempo di silenzio $t_{off}=50 \mu sec$. Ogni burst ha una durata $t_{on}$ che puo' variare da $0.1 sec$ a $4 sec$. In ciascun periodo $t_{on}$ arrivano sempre 2MB, con arrivi distribuiti uniformente. Assumendo una gestione con polling, calcolare il valore max di $T_{TR}$, assumendo $W=0.1$ e $X_M = 100$, affinche' possa essere smaltito correttamente il traffico bursty nell'ipotesi di $t_{ON\_min}=1 sec$.}

\begin{enumerate}[label=\arabic*.]
\item $3.5 \mu sec$
\item Nessuna delle altre opzioni
\item $0.045\times10^{-6}$ sec/byte
\item $0.005\times10^{-8}$ sec/byte
\end{enumerate}
\textbf{Q16: In un sistema di gestione dell’I/O con DMA si hanno migliori prestazioni per il trasferimento di N byte rispetto ad un sistema di gestione basato sull’interrupt soltanto se:}

\begin{enumerate}[label=\arabic*.]
\item Si utilizza una architettura Harvard
\item Non ci sono conflitti (f = 0)
\item N è sufficientemente alto
\item Nessuna delle altre opzioni
\end{enumerate}
\textbf{Q17: Data una gestione dell’I/O basata su DMA per trasferire burst di dati dall’I/O e che un periodo di silenzio $T_{off}$ di ciascun burst venga utilizzato per programmare il DMA e per gestire l’interrupt di fine trasferimento. Qual è il massimo throughput che si può smaltire in ciascun burst, nel periodo $T_{ON}$ in cui arrivano N bytes, nel caso di una architettura Harvard, percentuale di conflitti f non nulla, e una politica di gestione dei conflitti di tipo "DMA waiting for CPU" ?}

\begin{enumerate}[label=\arabic*.]
\item Nessuna delle altre opzioni
\item $\frac{N}{s+ 2 \times f \times N + N} \times \frac{1}{T_{tr\_DMA}} $
\item $\frac{1}{T_{tr\_DMA}}$
\item $0.5 \times \frac{N}{s+ 2 \times f \times N + N} \times \frac{1}{T_{tr\_DMA}} $
\end{enumerate}
\textbf{Q18: In un sistema IoT in cui sono presenti 16 sorgenti analogiche identiche, ciascuna caratterizzata da una $f_{MAX[segnale]}$ pari a K. Si decide di utilizzare un multiplexer analogico 16:1 che alimenta a valle un Sample\\\&Hold e un convertitore A/D nella catena di acquisizione. Selezionare quale delle seguenti opzioni è vera:}

\begin{enumerate}[label=\arabic*.]
\item Nessuna delle altre opzioni
\item Il componente Sample\&Hold a valle può mantenere la stessa performance richiesta nel caso vi fosse un'unica sorgente con $f_{MAX[segnale]} = K$
\item Il componente Sample\&Hold a valle può mantenere la stessa performance richiesta nel caso vi fosse un'unica sorgente con $f_{MAX[segnale]} = K$, purché il convertitore a valle sia 32 volte più performante
\item Il componente Sample\&Hold a valle deve avere una performance almeno pari a 32 volte quella richiesta nel caso vi fosse un'unica sorgente con $f_{MAX[segnale]} = K$, purché il convertitore a valle abbia anch’esso una performance 32 volte superiore
\end{enumerate}
\textbf{Q19: In un sistema di I/O con interrupt il massimo throughput di I/O ottenibile scambiando N byte}

\begin{enumerate}[label=\arabic*.]
\item Nessuna delle precedenti
\item Dipende da N
\item Non dipende da N
\item Non dipende dall’interrupt overhead W
\end{enumerate}
\textbf{Q20: In un sistema di I/O gestito con il polling, la massima frequenza di segnale $f_{MAX[segnale]}$ acquisibile:}

\begin{enumerate}[label=\arabic*.]
\item Nessuna delle altre opzioni
\item E' direttamente proporzionale al valore $T_{TR}$
\item Non dipende da $X_M$
\item Non dipende da W
\end{enumerate}
\textbf{Q21: Con riferimento all’architettura generale di un sistema IoT dire quale fra le seguenti opzioni è quella vera}

\begin{enumerate}[label=\arabic*.]
\item Il Fog Computing serve soprattutto a interfacciare i nodi IoT con la piattaforma IoT
\item Il Fog Computing non si utilizza nel caso vi siano vincoli troppo stringenti di real time
\item il Fog Computing serve soprattutto a esporre interfacce di programmazione più agevoli agli sviluppatori di applicazioni, nascondendo loro la complessità dei livelli più bassi della gerarchia
\item Nessuna delle altre risposte
\end{enumerate}
\end{multicols}

\newpage

\section*{IoT Systems and Technologies \footnotesize [B99CB8CB5]}
{\footnotesize Nome:\_\_\_\_\_\_\_\_\_\_\_\_\_\_\_\_ Cognome:\_\_\_\_\_\_\_\_\_\_\_\_\_\_\_\_ Mat:\_\_\_\_\_\_\_\_}
\begin{multicols}{2}

\textbf{Q1: Cosa contiene la memoria flash di un microcontrollore?}

\begin{enumerate}[label=\arabic*.]
\item I dati per il debug mediante connessione seriale
\item I dati temporanei delle variabili
\item Le connessioni di rete attive
\item Il programma da eseguire
\end{enumerate}
\textbf{Q2: Il Successive Approximation Converter:}

\begin{enumerate}[label=\arabic*.]
\item Risolve i problemi dovuti ai pin floating senza pull resistor
\item Risolve i problemi del Binary Weighted Resistor converter
\item Risolve i problemi di meta-stabilità
\item Risolve i problemi del Tracking Converter
\end{enumerate}
\textbf{Q3: Un possibile difetto del convertitore A/D mediante Tracking:}

\begin{enumerate}[label=\arabic*.]
\item Può portare alla meta-stabilità
\item Tutte le affermazioni sono corrette
\item Fa uso della ricerca binaria, con conseguente imprecisione
\item Potrebbe essere lento nell'inseguire il segnale da convertire
\end{enumerate}
\textbf{Q4: Qual è la differenza tra full-duplex e half-duplex?}

\begin{enumerate}[label=\arabic*.]
\item Half-duplex richiede più memoria
\item Full-duplex è sempre asincrono
\item Nel full-duplex entrambi possono trasmettere contemporaneamente
\item Nel full-duplex si usa una sola linea dati
\end{enumerate}
\textbf{Q5: Quali sono le quattro linee dell’interfaccia SPI?}

\begin{enumerate}[label=\arabic*.]
\item IN, OUT, GND, VCC
\item TXD, RXD, CLK, EN
\item MOSI, MISO, SCK, SS
\item SDA, SCL, ACK, RST
\end{enumerate}
\textbf{Q6: Quali sono i pin principali di una UART?}

\begin{enumerate}[label=\arabic*.]
\item SDA e SCL
\item MOSI e MISO
\item CLK e EN
\item TX e RX
\end{enumerate}
\textbf{Q7: Per un LED, tipicamente è necessaria:}

\begin{enumerate}[label=\arabic*.]
\item Resistenza di qualche centinaio di Ohm in parallelo
\item Resistenza di qualche Mega Ohm in parallelo
\item Resistenza di qualche Mega Ohm in serie
\item Resistenza di qualche centinaio di Ohm in serie
\end{enumerate}
\textbf{Q8: L'upload del programma alla memoria flash del microcontrollore è necessario:}

\begin{enumerate}[label=\arabic*.]
\item Ogni qual volta si modifica il circuito sulla breadboard
\item Ogni qual volta si esegue il codice del programma
\item Ogni qual volta si toglie l'alimentazione alla board
\item Ogni qual volta si modifica il codice del programma
\end{enumerate}
\textbf{Q9: Perché è sconsigliato usare `delay()` in applicazioni complesse?}

\begin{enumerate}[label=\arabic*.]
\item Blocca l’esecuzione del programma
\item Disattiva gli interrupt
\item Richiede molta memoria
\item Consuma più energia
\end{enumerate}
\textbf{Q10: Per quale motivo è utile il parametro $INPUT\_PULLUP$?}

\begin{enumerate}[label=\arabic*.]
\item Per aumentare la risoluzione PWM
\item Per attivare una resistenza interna
\item Per collegare un buzzer
\item Per leggere segnali analogici
\end{enumerate}
\textbf{Q11: Qual è l’uso principale della funzione `attachInterrupt()`?}

\begin{enumerate}[label=\arabic*.]
\item Gestire eventi su un GPIO con una ISR
\item Attivare la comunicazione SPI
\item Scrivere dati su un pin digitale
\item Controllare segnali PWM
\end{enumerate}
\textbf{Q12: A cose serve la funzione $Blink.run()$}

\begin{enumerate}[label=\arabic*.]
\item Apre una nuova connessione Wi-Fi per gli interrupt
\item Gestisce gli interrupt in alternativa ad attachInterrupt()
\item Configura i pin GPIO per il cloud dentro setup()
\item Gestisce la connessione cloud e messaggi virtuali dentro loop()
\end{enumerate}
\textbf{Q13: Dato una gestione dell’I/O basata su DMA per trasferire burst di dati, si assuma che il periodo di silenzio  ($T_{off}$) di ciascun burst venga utilizzato per programmare il DMA e per gestire l’interrupt di fine trasferimento e che nel periodo $T_{ON} = 0.2$ sec arrivino un massimo di $N=10^6$ bytes. Calcolare il numero massimo $N_{max}$ di bytes in ciascun burst che possono essere trasferiti attraverso il DMA in una architettura Harvard con percentuale di conflitti $f = 0.3$ gestiti con DMA waits for CPU e dato $T_{LOAD\/STORE} = 2 \times T_{tr\_DMA}$ e $T_{tr\_DMA} = 10^{-6}$ sec/byte.}

\begin{enumerate}[label=\arabic*.]
\item $0.25 \times 10^5$ bytes
\item $0.71 \times 2 \times 10^6$ bytes
\item $2 \times 10^6$ bytes
\item $0.125 \times 10^6$ bytes
\end{enumerate}
\textbf{Q14: In un sistema IoT occorre gestire l'acquisizione di burst di dati caratterizzati da un tempo di silenzio $t_{off}=50\mu sec$. Ogni burst ha una durata $t_{ON}$ che puo' variare da $t_{ON\_min}$ uguale a $0.1 sec$ a $4 sec$. In ciascun periodo $t_{ON}$ arrivano sempre 2MB, con arrivi distribuiti uniformemente.  Assumendo una gestione con interrupt e $T_{ISR}=1 \mu sec$ calcolare il $W_{MAX}$ (interrupt overhead) per smaltire correttamente il traffico bursty sopra descritto.  Indicare l'opzione corretta:}

\begin{enumerate}[label=\arabic*.]
\item 0.1
\item 0.5
\item Nessuna delle altre opzioni
\item 0.05
\end{enumerate}
\textbf{Q15: Assumendo una gestion del sistema di I/O con polling, con $T_{TR}=100\mu sec$, calcolare il minimo valore di $\eta_{Poll}$ affinche' si possano acquisire 5 sorgenti analogiche ciascuna caratterizzata da una $f_{MAX[segnale]}$ di 0.3Kbyte/sec.  Indicare quale dei seguenti valore è quello corretto:}

\begin{enumerate}[label=\arabic*.]
\item 0.3
\item 0.27
\item 0.15
\item 0.35
\end{enumerate}
\textbf{Q16: Data una gestione dell’I/O basata su DMA per trasferire burst di dati dall’I/O e che un periodo di silenzio $T_{off}$ di ciascun burst venga utilizzato per programmare il DMA e per gestire l’interrupt di fine trasferimento. Qual è il massimo throughput che si può smaltire in ciascun burst, nel periodo $T_{ON}$ in cui arrivano N bytes, nel caso di una architettura Harvard, percentuale di conflitti f non nulla, e una politica di gestione dei conflitti di tipo "DMA waiting for CPU" e assumendo $T_{LOAD\_STORE} = 2 \times T_{tr\_DMA}$}

\begin{enumerate}[label=\arabic*.]
\item $\frac{1}{ 2 \times f + 1} \times \frac{1}{T_{tr\_DMA}} $
\item $\frac{1}{T_{tr\_DMA}}$
\item $\frac{N}{s+ 2 \times f \times N + N} \times \frac{1}{T_{tr\_DMA}} $
\item $0.5 \times \frac{N}{s+ 2 \times f \times N + N} \times \frac{1}{T_{tr\_DMA}} $
\end{enumerate}
\textbf{Q17: Si utilizzi un sistema di gestione dell’I/O basato sul DMA per trasferire N bytes di dati dall’I/O. Si assuma una architettura Harvard, una percentuale di conflitti f diversa da 0 e una politica "CPU waits for DMA". Quale fra le seguenti affermazioni è corretta:}

\begin{enumerate}[label=\arabic*.]
\item Il $CPU_{TIME\_IO}$ che la CPU spende per l’I/O è trascurabile per valori di N molto grandi
\item Il $CPU_{TIME\_IO}$ che la CPU spende per l’I/O non dipende da f
\item Nessuna delle altre opzioni
\item Il $CPU_{TIME\_IO}$ che la CPU spende per l’I/O dipende dal prodotto $f \times N$
\end{enumerate}
\textbf{Q18: In un sistema IoT con due sorgenti analogiche caratterizzate dallo stesso valore di $f_{MAX[segnale]}$ espresso in byte/sec, assumendo una catena di acquisizione con un unico convertitore A/D, selezionare quale delle seguenti opzioni è vera:}

\begin{enumerate}[label=\arabic*.]
\item Il tempo massimo di conversione del convertitore deve essere minore o uguale a $(\frac{1}{4 \times f_{MAX[segnale]}})$
\item Il tempo massimo di conversione del convertitore deve essere minore o uguale a $(\frac{1}{f_{MAX[segnale]}})$
\item Nessuna delle opzioni
\item Il tempo massimo di conversione del convertitore deve essere minore o uguale a $(\frac{1}{2 \times f_{MAX[segnale]}})$
\end{enumerate}
\textbf{Q19: In un sistema di I/O con interrupt il massimo throughput di I/O ottenibile scambiando N byte}

\begin{enumerate}[label=\arabic*.]
\item Non dipende dall’interrupt overhead W
\item Non dipende da N
\item Nessuna delle precedenti
\item Dipende da N
\end{enumerate}
\textbf{Q20: In un sistema di I/O gestito con il polling, la massima frequenza di segnale $f_{MAX[segnale]}$ acquisibile:}

\begin{enumerate}[label=\arabic*.]
\item Non dipende da $X_M$
\item Nessuna delle altre opzioni
\item Non dipende da W
\item E' direttamente proporzionale al valore $T_{TR}$
\end{enumerate}
\textbf{Q21: Con riferimento all’architettura generale di un sistema IoT dire quale fra le seguenti opzioni è quella vera}

\begin{enumerate}[label=\arabic*.]
\item Il Fog Computing non si utilizza nel caso vi siano vincoli troppo stringenti di real time
\item il Fog Computing serve soprattutto a esporre interfacce di programmazione più agevoli agli sviluppatori di applicazioni, nascondendo loro la complessità dei livelli più bassi della gerarchia
\item Nessuna delle altre risposte
\item Il Fog Computing serve soprattutto a interfacciare i nodi IoT con la piattaforma IoT
\end{enumerate}
\end{multicols}

\newpage

\section*{IoT Systems and Technologies \footnotesize [B99CB8CB6]}
{\footnotesize Nome:\_\_\_\_\_\_\_\_\_\_\_\_\_\_\_\_ Cognome:\_\_\_\_\_\_\_\_\_\_\_\_\_\_\_\_ Mat:\_\_\_\_\_\_\_\_}
\begin{multicols}{2}

\textbf{Q1: Quale memoria conserva i dati anche se il dispositivo è spento?}

\begin{enumerate}[label=\arabic*.]
\item DRAM (Dynamic RAM)
\item Stack
\item SRAM (Static RAM)
\item Flash
\end{enumerate}
\textbf{Q2: In un microcontrollore, l'output PWM di un pin può essere utilizzato}

\begin{enumerate}[label=\arabic*.]
\item Per estendere il range di valori di un timer convenzionale
\item Per generare un segnale di clock interno a risoluzione maggiore
\item Per programmare le Interrupt Service Routine
\item Per implementare un convertitore da digitale ad analogico
\end{enumerate}
\textbf{Q3: Come funziona il convertitore SAR (Successive Approximation)?}

\begin{enumerate}[label=\arabic*.]
\item Usa una memoria cache
\item Conta fino al valore massimo
\item Sincronizza con il clock
\item Usa una ricerca binaria
\end{enumerate}
\textbf{Q4: In una interfaccia single-ended}

\begin{enumerate}[label=\arabic*.]
\item Il trasmettitore ed il ricevitore fanno riferimento alla stessa massa
\item Ogni bit è trasmesso su due linee singole (cioè non comunicanti tra loro)
\item Ogni bit trasmesso è seguito da un bit di end
\item E' meno soggetta ad interferenze e rumore
\end{enumerate}
\textbf{Q5: L’interfaccia I2C è...}

\begin{enumerate}[label=\arabic*.]
\item Asincrona, point-to-point
\item Solo per segnali analogici
\item Full-duplex, master-master
\item Sincrona, bus-based, half-duplex
\end{enumerate}
\textbf{Q6: L'interfaccia SPI è...}

\begin{enumerate}[label=\arabic*.]
\item Asincrona, broadcast
\item Sincrona, punto a punto, master-slave
\item Differenziale e half-duplex
\item Parallela e half-duplex
\end{enumerate}
\textbf{Q7: A cosa serve la fase di calibrazione nei sensori luce?}

\begin{enumerate}[label=\arabic*.]
\item Per impostare la frequenza del LED
\item Per impostare il colore del LED
\item Per rilevare il range di valori utili
\item Per stabilire un valore fisso di tensione
\end{enumerate}
\textbf{Q8: La calibrazione di un sensore potrebbe essere necessaria:}

\begin{enumerate}[label=\arabic*.]
\item Per determinare la corrente massima assorbita
\item Per determinare il range di variazione delle tensioni lette
\item Per riavviare il circuito del sensore in caso di inutilizzo prolungato
\item Per evitare che la DigitalRead restituisca valori intermedi tra 0 ed 1
\end{enumerate}
\textbf{Q9: Quando si usa dacWrite invece di PWM?}

\begin{enumerate}[label=\arabic*.]
\item Quando si vuole risparmiare memoria
\item Quando serve un segnale digitale
\item Quando serve un’uscita analogica precisa
\item Quando si usano motori
\end{enumerate}
\textbf{Q10: Su quali pin è possibile utilizzare `dacWrite()`?}

\begin{enumerate}[label=\arabic*.]
\item Su tutti i pin analogici
\item Su tutti i pin digitali
\item Su pin che vanno scelti in fase di setup
\item Solo su alcuni pin
\end{enumerate}
\textbf{Q11: Cosa è inopportuno porre dentro una ISR:}

\begin{enumerate}[label=\arabic*.]
\item Un insieme di istruzioni che modificano variabili locali
\item Un insieme di istruzioni che usi variabili booleane
\item Un insieme di istruzioni ad alta latenza
\item Un insieme di istruzioni che modificano variabili globali
\end{enumerate}
\textbf{Q12: A cose serve la funzione $Blink.run()$}

\begin{enumerate}[label=\arabic*.]
\item Configura i pin GPIO per il cloud dentro setup()
\item Gestisce gli interrupt in alternativa ad attachInterrupt()
\item Gestisce la connessione cloud e messaggi virtuali dentro loop()
\item Apre una nuova connessione Wi-Fi per gli interrupt
\end{enumerate}
\textbf{Q13: Dato una gestione dell’I/O basata su DMA per trasferire burst di dati, si assuma che il periodo di silenzio  ($T_{off}$) di ciascun burst venga utilizzato per programmare il DMA e per gestire l’interrupt di fine trasferimento e che nel periodo $T_{ON} = 0.2$ sec arrivino un massimo di $N=10^6$ bytes. Calcolare il numero massimo $N_{max}$ di bytes in ciascun burst che possono essere trasferiti attraverso il DMA in una architettura Harvard con percentuale di conflitti $f = 0.3$ gestiti con DMA waits for CPU e dato $T_{LOAD\/STORE} = 2 \times T_{tr\_DMA}$ e $T_{tr\_DMA} = 10^{-6}$ sec/byte.}

\begin{enumerate}[label=\arabic*.]
\item $0.25 \times 10^5$ bytes
\item $2 \times 10^6$ bytes
\item $0.125 \times 10^6$ bytes
\item $0.71 \times 2 \times 10^6$ bytes
\end{enumerate}
\textbf{Q14: Sia dato un sistema IoT in cui si voglia valutare la massima frequenza di segnale acquisibile da una sorgente analogica, espressa in byte/sec. Assumendo una gestione del sistema di I/O con Interrupt, con $T_{ISR}=100\mu sec$ e $W=0.1$, indicare la risposta corretta:}

\begin{enumerate}[label=\arabic*.]
\item $0.32\times10^3 byte/sec$
\item $0.45\times10^4 byte/sec$
\item $0.22\times10^5 byte/sec$
\item $0.045\times10^3 byte/sec$
\end{enumerate}
\textbf{Q15: In un sistema IoT occorre gestire l'acquisizione di un burst di dati con tempo di silenzio $t_{off}=50 \mu sec$. Ogni burst ha una durata $t_{on}$ che puo' variare da $0.1 sec$ a $4 sec$. In ciascun periodo $t_{on}$ arrivano sempre 2MB, con arrivi distribuiti uniformente. Assumendo una gestione con polling, calcolare il valore max di $T_{TR}$, assumendo $W=0.1$ e $X_M = 100$, affinche' possa essere smaltito correttamente il traffico bursty nell'ipotesi di $t_{ON\_min}=1 sec$.}

\begin{enumerate}[label=\arabic*.]
\item Nessuna delle altre opzioni
\item $0.005\times10^{-8}$ sec/byte
\item $0.045\times10^{-6}$ sec/byte
\item $3.5 \mu sec$
\end{enumerate}
\textbf{Q16: In un sistema di I/O gestito con il DMA, architettura Harvard, $f=0.25$, il max throughput di I/O ottenibile scambiando N bytes:}

\begin{enumerate}[label=\arabic*.]
\item E' almeno il doppio di quello ottenibile con $f=0$
\item Nessuna delle altre opzioni
\item E' inversamente proporzionale al valore N
\item E' direttamente proporzionale al valore N
\end{enumerate}
\textbf{Q17: Data una gestione dell’I/O basata su DMA per trasferire burst di dati dall’I/O e che un periodo di silenzio $T_{off}$ di ciascun burst venga utilizzato per programmare il DMA e per gestire l’interrupt di fine trasferimento. Qual è il massimo throughput che si può smaltire in ciascun burst, nel periodo $T_{ON}$ in cui arrivano N bytes, nel caso di una architettura Harvard, percentuale di conflitti f non nulla, e una politica di gestione dei conflitti di tipo "DMA waiting for CPU" e assumendo $T_{LOAD\_STORE} = 2 \times T_{tr\_DMA}$}

\begin{enumerate}[label=\arabic*.]
\item $0.5 \times \frac{N}{s+ 2 \times f \times N + N} \times \frac{1}{T_{tr\_DMA}} $
\item $\frac{1}{ 2 \times f + 1} \times \frac{1}{T_{tr\_DMA}} $
\item $\frac{N}{s+ 2 \times f \times N + N} \times \frac{1}{T_{tr\_DMA}} $
\item $\frac{1}{T_{tr\_DMA}}$
\end{enumerate}
\textbf{Q18: In un sistema IoT con due sorgenti analogiche caratterizzate dallo stesso valore di $f_{MAX[segnale]}$ espresso in byte/sec, assumendo una catena di acquisizione con un unico convertitore A/D, selezionare quale delle seguenti opzioni è vera:}

\begin{enumerate}[label=\arabic*.]
\item Nessuna delle opzioni
\item Il massimo throughput di I/O deve essere minore o uguale a $2 \times f_{MAX[segnale]}$ byte/sec
\item Il massimo throughput di I/O deve essere minore o uguale a $8 \times f_{MAX[segnale]}$ byte/sec
\item Il massimo throughput di I/O deve essere almeno pari a $4 \times f_{MAX[segnale]}$ byte/sec
\end{enumerate}
\textbf{Q19: Considerando l’utilizzo dell’interrupt come tecnica di gestione dell’I/O:}

\begin{enumerate}[label=\arabic*.]
\item La gestione della priorità serve ad abbattere l’interrupt overhead
\item L’interrupt overhead W può essere diminuito se si utilizza una gestione vettorizzata dell’interrupt
\item Il vettore di interrupt serve principalmente ad indicare alla CPU la priorità di gestione della ISR
\item Nessuna delle opzioni riportate
\end{enumerate}
\textbf{Q20: In un sistema di I/O gestito con il polling, la massima frequenza di segnale $f_{MAX[segnale]}$ acquisibile:}

\begin{enumerate}[label=\arabic*.]
\item Non dipende da W
\item E' inversamente proporzionale al valore $T_{TR}$
\item Nessuna delle altre opzioni
\item Non dipende da $X_M$
\end{enumerate}
\textbf{Q21: Con riferimento all’architettura generale di un sistema IoT dire quale fra le seguenti opzioni è quella vera}

\begin{enumerate}[label=\arabic*.]
\item L’uso congiunto di Fog e MIST computing in un sistema IoT normalmente serve ad evitare l’utilizzo del Cloud Computing
\item Nessuna delle altre risposte
\item Il ruolo del MIST Computing è principalmente quello di ridurre la complessità del software a bordo dei nodi IoT per limitare il consumo di potenza del nodo
\item L’uso del MIST Computing, in alcuni scenari che prevedono comunicazioni wireless, comporta un maggior consumo di energia del nodo IoT dovuto alla trasmissione dati
\end{enumerate}
\end{multicols}

\newpage

\section*{IoT Systems and Technologies \footnotesize [B99CB8CB7]}
{\footnotesize Nome:\_\_\_\_\_\_\_\_\_\_\_\_\_\_\_\_ Cognome:\_\_\_\_\_\_\_\_\_\_\_\_\_\_\_\_ Mat:\_\_\_\_\_\_\_\_}
\begin{multicols}{2}

\textbf{Q1: Cosa contiene la memoria SRAM di un microcontrollore?}

\begin{enumerate}[label=\arabic*.]
\item Le variabili in uso dal programma
\item Il sistema operativo
\item La configurazione del microcontrollore
\item I dati di debug
\end{enumerate}
\textbf{Q2: Il convertitore A/D mediante Tracking:}

\begin{enumerate}[label=\arabic*.]
\item Utilizza internamente un convertitore digitale/analogico
\item Necessita di un trigger di Schmidt
\item Insegue il valore attuale di tensione in input al fine di stabilizzarlo
\item Utilizza internamente un convertitore analogico/digitale
\end{enumerate}
\textbf{Q3: La risoluzione del convertitore A/D è di 10 bit per un range di tensioni da 0 a 5V. L'acquisizione del valore 1023 identifica}

\begin{enumerate}[label=\arabic*.]
\item Un valore di tensione di circa 1.5V
\item Un valore di tensione fuori range
\item Un valore di tensione di circa 2.5V
\item Nessuna delle opzioni riportate
\end{enumerate}
\textbf{Q4: In una interfaccia seriale}

\begin{enumerate}[label=\arabic*.]
\item Richiede sempre una linea per il controllo di parità
\item I bit sono inviati in parallelo ma il modulo di destinazione li serializza
\item I bit sono inviati in sequenza un bit alla volta
\item Si utilizzano diverse linee fisiche al fine di trasmettere più bit contemporaneamente
\end{enumerate}
\textbf{Q5: Cosa fa un dispositivo I2C se ha bisogno di più tempo?}

\begin{enumerate}[label=\arabic*.]
\item Aumenta il baud rate
\item Blocca la linea SDA
\item Invia un bit di pausa
\item Tiene la linea di clock bassa
\end{enumerate}
\textbf{Q6: In una interfaccia I2C, il significato del valore che si trova nel segnale SDA:}

\begin{enumerate}[label=\arabic*.]
\item Dipende dal livello di tensione analogico che si trova su SDA
\item Dipende dalla frequenza di oscillazione del segnale, ogni frequenza avente un significato diverso
\item Serve a selezionare lo slave con cui il master vuole comunicare
\item Dipende dal valore attuale del segnale di clock SCL
\end{enumerate}
\textbf{Q7: Se connetto una board ESP32 all'alimentazione:}

\begin{enumerate}[label=\arabic*.]
\item Nessuna delle altre risposte
\item Non esegue niente finchè non si immette un nuovo codice sulla flash
\item Inizia ad eseguire l'ultimo codice caricato sulla flash
\item Azzera il proprio contenuto flash
\end{enumerate}
\textbf{Q8: Qual è il problema del “floating pin” quando un pulsante è aperto?}

\begin{enumerate}[label=\arabic*.]
\item Il pin può assumere valori casuali
\item Il LED si spegne automaticamente
\item Il pin si danneggia
\item La corrente diventa negativa
\end{enumerate}
\textbf{Q9: Quando si usa dacWrite invece di PWM?}

\begin{enumerate}[label=\arabic*.]
\item Quando serve un’uscita analogica precisa
\item Quando si usano motori
\item Quando si vuole risparmiare memoria
\item Quando serve un segnale digitale
\end{enumerate}
\textbf{Q10: Quale tra questi non è un metodo per generare un’uscita analogica su ESP32?}

\begin{enumerate}[label=\arabic*.]
\item analogWrite()
\item ledcWrite()
\item dacWrite()
\item Serial.write()
\end{enumerate}
\textbf{Q11: Impostando CHANGE come condizione interrupt su pin 7, avviene che:}

\begin{enumerate}[label=\arabic*.]
\item La ISR è invocata su di un fronte di salita o discesa nel pin 7
\item Il pin 7 è letto dalla ISR sul fronte di salita o discesa del clock
\item La ISR è invocata quando si modifica il numero di interrupt associato al pin 7
\item Tutte le ISR del codice sono invocate quando vi è un cambiamento del pin 7
\end{enumerate}
\textbf{Q12: Cosa rappresenta un "Device Template" in Blynk?}

\begin{enumerate}[label=\arabic*.]
\item Un'app per smartphone per progettare interfacce IoT
\item Un modello per progettare sensori
\item Un modello per creare dispositivi con le stesse funzionalità
\item Un firmware aggiornato per modelli con simili funzionalità
\end{enumerate}
\textbf{Q13: Dato una gestione dell’I/O basata su DMA per trasferire burst di dati, si assuma che il periodo di silenzio  ($T_{off}$) di ciascun burst venga utilizzato per programmare il DMA e per gestire l’interrupt di fine trasferimento e che nel periodo $T_{ON} = 0.2$ sec arrivino un massimo di $N=10^6$ bytes. Calcolare il numero massimo $N_{max}$ di bytes in ciascun burst che possono essere trasferiti attraverso il DMA in una architettura Harvard con percentuale di conflitti $f = 0.3$ gestiti con DMA waits for CPU e dato $T_{LOAD\/STORE} = 2 \times T_{tr\_DMA}$ e $T_{tr\_DMA} = 10^{-6}$ sec/byte.}

\begin{enumerate}[label=\arabic*.]
\item $0.125 \times 10^6$ bytes
\item $2 \times 10^6$ bytes
\item $0.25 \times 10^5$ bytes
\item $0.71 \times 2 \times 10^6$ bytes
\end{enumerate}
\textbf{Q14: In un sistema IoT occorre gestire l'acquisizione di burst di dati caratterizzati da un tempo di silenzio $t_{off}=50\mu sec$. Ogni burst ha una durata $t_{ON}$ che puo' variare da $t_{ON\_min}$ uguale a $0.1 sec$ a $4 sec$. In ciascun periodo $t_{ON}$ arrivano sempre 2MB, con arrivi distribuiti uniformemente.  Assumendo una gestione con interrupt e $T_{ISR}=1 \mu sec$ calcolare il $W_{MAX}$ (interrupt overhead) per smaltire correttamente il traffico bursty sopra descritto.  Indicare l'opzione corretta:}

\begin{enumerate}[label=\arabic*.]
\item 0.05
\item Nessuna delle altre opzioni
\item 0.5
\item 0.1
\end{enumerate}
\textbf{Q15: Assumendo una gestion del sistema di I/O con polling, con $T_{TR}=100\mu sec$, calcolare il minimo valore di $\eta_{Poll}$ affinche' si possano acquisire 5 sorgenti analogiche ciascuna caratterizzata da una $f_{MAX[segnale]}$ di 0.3Kbyte/sec.  Indicare quale dei seguenti valore è quello corretto:}

\begin{enumerate}[label=\arabic*.]
\item 0.15
\item 0.3
\item 0.35
\item 0.27
\end{enumerate}
\textbf{Q16: Si utilizzi un sistema di gestione dell’I/O basato sul DMA per trasferire N bytes di dati dall’I/O. Si assuma una architettura Harvard, una percentuale di conflitti f diversa da 0 e una politica "CPU waits for DMA". Quale fra le seguenti affermazioni è corretta:}

\begin{enumerate}[label=\arabic*.]
\item Nessuna delle altre opzioni
\item Il $CPU_{TIME\_IO}$ che la CPU spende per l’I/O non dipende da N
\item Il $CPU_{TIME\_IO}$ che la CPU spende per l’I/O diminuisce al crescere di N
\item A causa del $CPU_{TIME\_IO}$ che la CPU spende per la gestione dei conflitti, superato un certo valore di f non si ha più alcun vantaggio ad utilizzare il DMA
\end{enumerate}
\textbf{Q17: Data una gestione dell’I/O basata su DMA per trasferire burst di dati dall’I/O e che un periodo di silenzio $T_{off}$ di ciascun burst venga utilizzato per programmare il DMA e per gestire l’interrupt di fine trasferimento. Qual è il massimo throughput che si può smaltire in ciascun burst, nel periodo $T_{ON}$ in cui arrivano N bytes, nel caso di una architettura Harvard, percentuale di conflitti f non nulla, e una politica di gestione dei conflitti di tipo "DMA waiting for CPU" e assumendo $T_{LOAD\_STORE} = 2 \times T_{tr\_DMA}$}

\begin{enumerate}[label=\arabic*.]
\item Nessuna delle altre opzioni
\item $\frac{1}{T_{tr\_DMA}}$
\item $0.5 \times \frac{N}{s+ 2 \times f \times N + N} \times \frac{1}{T_{tr\_DMA}} $
\item $\frac{N}{s+ 2 \times f \times N + N} \times \frac{1}{T_{tr\_DMA}} $
\end{enumerate}
\textbf{Q18: In un sistema IoT con due sorgenti analogiche caratterizzate dallo stesso valore di $f_{MAX[segnale]}$ espresso in byte/sec, assumendo una catena di acquisizione con un unico convertitore A/D, selezionare quale delle seguenti opzioni è vera:}

\begin{enumerate}[label=\arabic*.]
\item Il tempo massimo di conversione del convertitore deve essere minore o uguale a $(\frac{1}{2 \times f_{MAX[segnale]}})$
\item Il tempo massimo di conversione del convertitore deve essere minore o uguale a $(\frac{1}{f_{MAX[segnale]}})$
\item Il tempo massimo di conversione del convertitore deve essere minore o uguale a $(\frac{1}{4 \times f_{MAX[segnale]}})$
\item Nessuna delle opzioni
\end{enumerate}
\textbf{Q19: In un sistema di I/O con interrupt il massimo throughput di I/O ottenibile scambiando N byte}

\begin{enumerate}[label=\arabic*.]
\item Non dipende da N
\item Dipende da N
\item Nessuna delle precedenti
\item Non dipende dall’interrupt overhead W
\end{enumerate}
\textbf{Q20: In un sistema di I/O gestito con il polling, la massima frequenza di segnale $f_{MAX[segnale]}$ acquisibile:}

\begin{enumerate}[label=\arabic*.]
\item Non dipende da $X_M$
\item Non dipende da W
\item E' direttamente proporzionale al valore $T_{TR}$
\item Nessuna delle altre opzioni
\end{enumerate}
\textbf{Q21: Con riferimento all’architettura generale di un sistema IoT:}

\begin{enumerate}[label=\arabic*.]
\item In uno scenario di veicoli a guida autonoma il fog computing serve a limitare la latenza nello scambio di informazioni soggette a vincoli di real time stringenti
\item Nessuna delle altre risposte
\item Il fog computing è normalmente residente a bordo dei nodi IoT
\item In una architettura di un sistema IoT che utilizza il Fog Computing, i nodi IoT devono essere più complessi per potere gestire l’interazione con i nodi fog
\end{enumerate}
\end{multicols}

\newpage

\section*{IoT Systems and Technologies \footnotesize [B99CB8CB8]}
{\footnotesize Nome:\_\_\_\_\_\_\_\_\_\_\_\_\_\_\_\_ Cognome:\_\_\_\_\_\_\_\_\_\_\_\_\_\_\_\_ Mat:\_\_\_\_\_\_\_\_}
\begin{multicols}{2}

\textbf{Q1: Cosa contiene la memoria SRAM di un microcontrollore?}

\begin{enumerate}[label=\arabic*.]
\item Le variabili in uso dal programma
\item I dati di debug
\item La configurazione del microcontrollore
\item Il sistema operativo
\end{enumerate}
\textbf{Q2: Qual è lo svantaggio del circuito resistivo binary-weight per DAC?}

\begin{enumerate}[label=\arabic*.]
\item Non funziona con segnali digitali
\item È troppo lento
\item Richiede resistenze di precisione
\item Ha alta latenza
\end{enumerate}
\textbf{Q3: Cosa rappresenta una risoluzione a 3 bit in un ADC?}

\begin{enumerate}[label=\arabic*.]
\item 3 livelli
\item 4 livelli
\item 8 livelli
\item 256 livelli
\end{enumerate}
\textbf{Q4: Qual è la caratteristica della comunicazione point-to-point?}

\begin{enumerate}[label=\arabic*.]
\item Usa un bus condiviso
\item Supporta broadcast
\item Avviene tra due soli dispositivi
\item Richiede indirizzamento complesso
\end{enumerate}
\textbf{Q5: In una interfaccia I2C:}

\begin{enumerate}[label=\arabic*.]
\item Seleziona uno slave mediante pilotando il segnale di clock
\item Il bus è costituito da due linee, più una terza linea per il clock
\item Il bus è costituito da due linee, di cui una è il clock
\item Presenta un bus con due linee delle MOSI e MISO
\end{enumerate}
\textbf{Q6: Come viene rilevato l’inizio di un pacchetto UART?}

\begin{enumerate}[label=\arabic*.]
\item Dall’assenza di segnale
\item Dal fronte di discesa del bit di start
\item Dal bit di parità
\item Dal clock condiviso
\end{enumerate}
\textbf{Q7: Qual è il problema del “floating pin” quando un pulsante è aperto?}

\begin{enumerate}[label=\arabic*.]
\item Il pin si danneggia
\item La corrente diventa negativa
\item Il LED si spegne automaticamente
\item Il pin può assumere valori casuali
\end{enumerate}
\textbf{Q8: Un sensore di tipo resistivo:}

\begin{enumerate}[label=\arabic*.]
\item Deve essere orientato in modo corretto
\item Deve essere messo in parallelo ad una resistenza
\item Deve essere messo in serie ad una resistenza
\item Ha 3 pin: alimentazione, massa e valore di tensione di uscita
\end{enumerate}
\textbf{Q9: Nel codice di ESP32 la funzione setup()}

\begin{enumerate}[label=\arabic*.]
\item Viene eseguita alla fine di ogni ripetizione della funzione loop()
\item Può essere eseguita in qualunque momento, ma una volta sola, mediante apposita chiamata
\item Viene eseguita una sola volta, all'uscita della funzione loop()
\item Viene eseguita una sola volta, ad ogni avvio del programma
\end{enumerate}
\textbf{Q10: Cosa indica `millis()` su ESP32?}

\begin{enumerate}[label=\arabic*.]
\item Pausa espressa in millisecondi
\item Millisecondi trascorsi dall'avvio
\item Durata di un ciclo PWM
\item Frequenza del clock in millisecondi
\end{enumerate}
\textbf{Q11: Perché il debounce è importante con gli interrupt?}

\begin{enumerate}[label=\arabic*.]
\item Per aumentare la luminosità del LED
\item Per risparmiare energia
\item Per evitare chiamate multiple della ISR
\item Per evitare la disconnessione del pin
\end{enumerate}
\textbf{Q12: Cosa rappresenta un "Device Template" in Blynk?}

\begin{enumerate}[label=\arabic*.]
\item Un firmware aggiornato per modelli con simili funzionalità
\item Un modello per creare dispositivi con le stesse funzionalità
\item Un modello per progettare sensori
\item Un'app per smartphone per progettare interfacce IoT
\end{enumerate}
\textbf{Q13: Dato una gestione dell’I/O basata su DMA per trasferire burst di dati, si assuma che il periodo di silenzio  ($T_{off}$) di ciascun burst venga utilizzato per programmare il DMA e per gestire l’interrupt di fine trasferimento e che nel periodo $T_{ON} = 0.2$ sec arrivino un massimo di $N=10^6$ bytes. Calcolare il numero massimo $N_{max}$ di bytes in ciascun burst che possono essere trasferiti attraverso il DMA in una architettura Harvard con percentuale di conflitti $f = 0.3$ gestiti con DMA waits for CPU e dato $T_{LOAD\/STORE} = 2 \times T_{tr\_DMA}$ e $T_{tr\_DMA} = 10^{-6}$ sec/byte.}

\begin{enumerate}[label=\arabic*.]
\item $2 \times 10^6$ bytes
\item $0.125 \times 10^6$ bytes
\item $0.25 \times 10^5$ bytes
\item $0.71 \times 2 \times 10^6$ bytes
\end{enumerate}
\textbf{Q14: Assumendo una gestione del sistema di I/O con Interrupt, con $T_{ISR}=100\mu sec$, calcolare il valore massimo di $W$ affinche' si possano acquisire 10 sorgenti analogiche ciascuna caratterizzata da una $f_{MAX[segnale]}$ di 0.3Kbyte/sec.}

\begin{enumerate}[label=\arabic*.]
\item 0.81
\item 0.5
\item 0.66
\item 0.73
\end{enumerate}
\textbf{Q15: In un sistema IoT occorre gestire l'acquisizione di un burst di dati con tempo di silenzio $t_{off}=50 \mu sec$. Ogni burst ha una durata $t_{on}$ che puo' variare da $0.1 sec$ a $4 sec$. In ciascun periodo $t_{on}$ arrivano sempre 2MB, con arrivi distribuiti uniformente. Assumendo una gestione con polling, calcolare il valore max di $T_{TR}$, assumendo $W=0.1$ e $X_M = 100$, affinche' possa essere smaltito correttamente il traffico bursty nell'ipotesi di $t_{ON\_min}=1 sec$.}

\begin{enumerate}[label=\arabic*.]
\item Nessuna delle altre opzioni
\item $0.005\times10^{-8}$ sec/byte
\item $0.045\times10^{-6}$ sec/byte
\item $3.5 \mu sec$
\end{enumerate}
\textbf{Q16: Si utilizzi un sistema di gestione dell’I/O basato sul DMA per trasferire N bytes di dati dall’I/O. Si assuma una architettura Harvard, una percentuale di conflitti f diversa da 0 e una politica "CPU waits for DMA". Quale fra le seguenti affermazioni è corretta:}

\begin{enumerate}[label=\arabic*.]
\item Il Throughput di I/O raddoppia se f si dimezza
\item Il Throughput di I/O dipende dal valore di f
\item Il Throughput di I/O diminuisce all’aumentare di f
\item Nessuna delle altre opzioni
\end{enumerate}
\textbf{Q17: Si utilizzi un sistema di gestione dell’I/O basato sul DMA per trasferire N bytes di dati dall’I/O. Si assuma una architettura Harvard, una percentuale di conflitti f diversa da 0 e una politica "CPU waits for DMA". Quale fra le seguenti affermazioni è corretta:}

\begin{enumerate}[label=\arabic*.]
\item Il $CPU_{TIME\_IO}$ che la CPU spende per l’I/O è trascurabile per valori di N molto grandi
\item Il $CPU_{TIME\_IO}$ che la CPU spende per l’I/O dipende dal prodotto $f \times N$
\item Il $CPU_{TIME\_IO}$ che la CPU spende per l’I/O non dipende da f
\item Nessuna delle altre opzioni
\end{enumerate}
\textbf{Q18: Con riferimento alla catena di acquisizione di un sistema IoT con una sorgente analogica caratterizzata da un valore di $f_{MAX[segnale]} = K$, selezionare quale delle seguenti opzioni è vera:}

\begin{enumerate}[label=\arabic*.]
\item Nessuna delle altre opzioni
\item La latenza del convertitore A/D deve essere almeno il doppio di $f_{MAX[segnale]}$
\item Il throughput minimo del sistema di I/O deve essere almeno il doppio della frequenza di campionamento di Nyquist
\item La latenza del Sample\&Hold deve essere il doppio di $f_{MAX[segnale]}$
\end{enumerate}
\textbf{Q19: Considerando una gestione dell’I/O con interrupt:}

\begin{enumerate}[label=\arabic*.]
\item Nessuna delle opzioni riportate
\item Il vettore di interrupt viene generato dalla CPU
\item L’utilizzo della tecnica dell’interrupt vettorizzato non diminuisce l’interrupt overhead W
\item Il vettore di interrupt elimina il problema di gestire la priorità
\end{enumerate}
\textbf{Q20: In un sistema di I/O gestito con il polling, il max throughput di I/O ottenibile:}

\begin{enumerate}[label=\arabic*.]
\item E' direttamente proporzionale al valore $T_{TR}$
\item Non dipende da W
\item Dipende da W
\item Nessuna delle altre opzioni
\end{enumerate}
\textbf{Q21: Con riferimento all’architettura generale di un sistema IoT dire quale fra le seguenti opzioni è quella vera}

\begin{enumerate}[label=\arabic*.]
\item Nessuna delle altre risposte
\item Il Fog Computing non si utilizza nel caso vi siano vincoli troppo stringenti di real time
\item il Fog Computing serve soprattutto a esporre interfacce di programmazione più agevoli agli sviluppatori di applicazioni, nascondendo loro la complessità dei livelli più bassi della gerarchia
\item Il Fog Computing serve soprattutto a interfacciare i nodi IoT con la piattaforma IoT
\end{enumerate}
\end{multicols}

\end{document}
